\section{Conclusion}

In this work, we have constructed a theory capable of formally describing classical Go from elementary notions—\emph{where, who, what, and how}—supported by set theory. Throughout this construction, the rules of the game emerge naturally, and intuitive concepts acquire a precise and tractable expression.

However, the principal value of this framework is not limited to formalizing Go. As explored in Part~\ref{parte3}, the developed structure is sufficiently general to encompass a broader family of games based on the idea of surrounding and capturing. This suggests the existence of a more abstract mathematical object, which we designate as \textsc{Surrounding Games} (\(\mathsf{SG}\)), of which Go is merely a particular case.

Several paths remain open: identifying which properties are exclusive to Go and which belong to the entire family, formalizing this family and seeing how far we can take it, constructing infinite Go, or attempting to connect this theory with combinatorial game theory.

\vfill

I'm actively thiking about this, I don't have my whole time for this but I will, from time to time make updates with new stuff. The thing I'm most intereted is attacking the infinte Go.